% Section 5.3, from lectures/L05/appendix/b_exercises.md.
%
% Exercise 11 is marked "(fördjupning)" after its kind rather than in its title; it gets the two
% kinds Förståelse and Fördjupning, like every other exercise so marked.

\section{Övningar}
\label{c5:sec:exercises}\label{c5:app:b}

\begin{aside}
\textbf{Så kontrollerar du ditt arbete.} De flesta övningarna här är felsökningsfall: en
beskrivning av ett nät, vad som mättes, och frågan var felet sitter. Svara alltid med \textbf{var du
mäter härnäst och varför}, inte bara med en gissning om orsaken. En gissning som råkar vara rätt
lär dig ingenting om nästa fel.

Kontrolluppgiften, övning~\ref{c5:ex:10}, görs på labbstationen, gärna under Labb~2-1.

Lösningsförslagen finns i bilaga~\ref{app:answers}.

Varje övning är märkt med sin sort.
\end{aside}

\exercise{Matningen först}{Förståelse}
\label{c5:ex:1}
\begin{parts}
  \item Mellan vilka ben mäter du matningen på en 74HC08, och vilket värde ska du se?
  \item Varför ska du mäta på kretsens ben och inte på matningsskenorna?
  \item En krets saknar matning helt, men nätet \enquote{fungerar nästan}. Hur kan en krets utan
    matning alls ge några utsignaler?
\end{parts}

\exercise{Vad betyder mätvärdet?}{Förståelse}
\label{c5:ex:2}
Du mäter på ingångar och utgångar i ett nät med 74HC-kretsar på 5~V. Vilken logisk nivå motsvarar
varje mätvärde, och vilka av dem tyder på ett fel?
\begin{parts}
  \item 4,7~V på en utgång.
  \item 0,1~V på en utgång.
  \item 2,3~V på en ingång.
  \item 2,5~V på en utgång, som dessutom är ansluten till en annan grinds utgång.
\end{parts}

\exercise{En grind som ljuger}{Räkna för hand}
\label{c5:ex:3}
På en krets som ska vara en 74HC08 mäter du på grind~1: ben~1 = 4,9~V, ben~2 = 0,1~V och ben~3 =
4,8~V.
\begin{parts}
  \item Vad borde ben~3 ha visat?
  \item Grinden beter sig som en annan grind. Vilken? Ge en trolig förklaring, och vad du
    kontrollerar för att bekräfta den.
\end{parts}

\exercise{Den varma kretsen}{Förståelse}
\label{c5:ex:4}
Direkt efter att du slagit på matningen känner du att en av kretsarna blir varm.
\begin{parts}
  \item Vad gör du först?
  \item Nämn två troliga orsaker.
  \item Hur kontrollerar du, utan att slå på matningen igen, vilken av dem det är?
\end{parts}

\exercise{Läs felet ur tabellen}{Räkna för hand}
\label{c5:ex:5}
Nätet \code{X = AB + C} ger den här tabellen:

\begin{narrowtable}{@{}cccc@{}}
  \thead{A} & \thead{B} & \thead{C} & \thead{X, mätt} \\
  \midrule
  0 & 0 & 0 & 0 \\
  0 & 0 & 1 & 0 \\
  0 & 1 & 0 & 0 \\
  0 & 1 & 1 & 0 \\
  1 & 0 & 0 & 0 \\
  1 & 0 & 1 & 0 \\
  1 & 1 & 0 & 1 \\
  1 & 1 & 1 & 1 \\
\end{narrowtable}

\begin{parts}
  \item Vilka rader är fel?
  \item Vad har de gemensamt, och vilken del av nätet pekar det ut?
  \item Var mäter du först?
\end{parts}

\exercise{Bromsassistenten som inte bromsar}{Räkna för hand}
\label{c5:ex:6}
Ett par har byggt bromsassistenten i \hyperref[lab2:1:u3]{Labb~2-1 uppgift~3}. Nätet stämmer på
tolv av sexton rader. De fyra felaktiga raderna är alla rader där \code{D = 1} och \code{F = 1}: där
blir \code{B = 0}, men ska vara 1. Beteckningarna är desamma som i \secref{c2:a10}.
\begin{parts}
  \item Vilket krav i specifikationen bryter nätet mot, och varför är just det felet farligt?
  \item Paret har kopplat förarens pedal på ett ställe i nätet där den inte hör hemma. Var sitter
    den, och vilket uttryck har de byggt i stället för det rätta?
  \item Var mäter du för att bekräfta det?
\end{parts}

\exercise{Halvering}{Förståelse}
\label{c5:ex:7}
En signal går genom sexton grindar i rad, och den är fel vid utgången. Den är rätt vid ingången.
\begin{parts}
  \item Var gör du den första mätningen, och varför just där?
  \item Hur många mätningar behöver du högst för att hitta den grind där felet uppstår?
  \item Hur många hade du behövt, i värsta fall, om du i stället börjat från ingången och mätt efter
    varje grind?
\end{parts}

\exercise{Den flimrande lysdioden}{Förståelse}
\label{c5:ex:8}
En lysdiod på utgången av ett nät tänds och släcks när du för handen nära kopplingsdäcket, fast du
inte rör strömbrytarna.
\begin{parts}
  \item Vad är den troliga orsaken?
  \item Hur hittar du exakt vilket ben det gäller?
  \item Varför händer det här med HC-kretsar, men inte på samma sätt med de äldre LS-kretsarna?
\end{parts}

\exercise{Planera fläkten}{Konstruktion}
\label{c5:ex:9}
Fläkten från övning~\ref{c4:ex:8} har uttrycket \code{F = TW'B' + CB'}.
\begin{parts}
  \item Gör en fullständig grindtilldelning med kretsbeteckningar och bennummer, med 74HC04,
    74HC08 och 74HC32.
  \item Skriv kopplingslistan, med matningen först och de oanvända ingångarna sist.
  \item Nätet fungerar, utom att fläkten aldrig startar på hög temperatur. Vilka rader i tabellen
    är fel, och vilken del av nätet mäter du först?
\end{parts}

\exercise{Kontroll: en utgång under last}{Kontroll}
\label{c5:ex:10}
\emph{Förutsäg, mät på labbstationen, förklara skillnaden.}

En utgång på en 74HC08 driver en röd lysdiod genom ett förkopplingsmotstånd på 330~Ω till jord.

\task{a) Förutsäg}
Räkna ut strömmen genom lysdioden som i \secref{c4:b5}, med antagandet att utgången ger exakt 5~V
och att lysdioden har 2~V över sig.

\task{b) Mät}
Med utgången hög: mät spänningen på utgången och spänningen över lysdioden. Med utgången låg: mät
spänningen på utgången.

\task{c) Räkna igen}
Nu med de uppmätta spänningarna: spänningen över motståndet är utgångens spänning minus
lysdiodens. Vilken ström går det?

\task{d) Förklara skillnaden}
Förklara skillnaden mellan a) och c). Vilket av de två antagandena i a) var fel, och varför?
Stämmer den uppmätta höga och låga nivån med databladets $V_{\text{OH}}$ och $V_{\text{OL}}$?

\exercise{Fel krets i sockeln}{Förståelse, Fördjupning}
\label{c5:ex:11}
Någon har satt en 74HC02 där det skulle sitta en 74HC00, och kopplat den enligt benplaceringen för
74HC00: strömbrytarna A och B till ben~1 och~2, och ben~3 till lysdioden.
\begin{parts}
  \item Vad är ben~1 på en 74HC02?
  \item Vad händer elektriskt när strömbrytaren på ben~1 står i läge 0 medan NOR-grinden försöker
    ge en etta på samma ben? Varför kan kretsen bli varm?
  \item Vad visar lysdioden på ben~3?
\end{parts}
